\section{The universal numerical upper bound}\label{sec:universal-bound}

The constant $0.792$ on the source page is not merely an asymptotic claim.
This section proves the exact all-positive-integer assertion of
Berend--Tassa~\cite{berend-tassa}, using a simple analytic estimate followed
by a small, completely rational certificate.  The estimate used here is
independent of the detailed saddle-point expansion.

\begin{theorem}[The $0.792$ bound]\label{thm:0792}
For every integer $n\ge1$,
\begin{equation}\label{eq:0792}
 B_n<\left(\frac{99n}{125\log(n+1)}\right)^n
     =\left(\frac{0.792n}{\log(n+1)}\right)^n.
\end{equation}
\end{theorem}
\begin{proof}
Take $r=\tfrac12\log(n+1)$ in
\eqref{eq:coefficient-bound}.  The elementary factorial estimate
\eqref{eq:elementary-factorial-log} gives
\begin{equation}\label{eq:0792-envelope}
 \frac{B_n^{1/n}\log(n+1)}n
 \le2\exp\left(-1+\frac{\log n+\sqrt{n+1}}n\right).
\end{equation}
The function $\log x/x$ decreases for $x\ge3$, since its derivative is
$(1-\log x)/x^2<0$.  The function $\sqrt{x+1}/x$ also decreases for $x>0$,
as direct differentiation gives
$-(x+2)/(2x^2\sqrt{x+1})<0$.
For $n\ge512$, therefore,
\[
 \frac{\log n+\sqrt{n+1}}n
 \le\frac{\log512+\sqrt{513}}{512}
 <\frac{7+23}{512}<\frac1{16}.
\]
Here $23^2>513$ and $e^7>512$ are elementary finite inequalities; for
example $e>5/2$ gives $(5/2)^7>512$.  Finally
\[
 e^{15/16}>\sum_{j=0}^{4}\frac{(15/16)^j}{j!}>\frac{250}{99}.
\]
The last comparison is a rational inequality verified by clearing
denominators.  Hence $2e^{-15/16}<99/125$, proving
\eqref{eq:0792} for every $n\ge512$.

The remaining finite assertion is
\begin{equation}\label{eq:small-bound-certificate}
 \bigwedge_{n=1}^{511}
 B_n\left(\frac{u_n}{2^{40}}\right)^n
       <\left(\frac{99n}{125}\right)^n,
 \qquad \log(n+1)\le\frac{u_n}{2^{40}}.
\end{equation}
We specify the integers $u_n$ by an exact bounded algorithm and supply all
511 values in \texttt{code/bound\_certificate.csv}.
Write $n+1=2^k y$ with $1\le y<2$, and $z=(y-1)/(y+1)$.  Define
\[
 L_J(z)=2\sum_{j=0}^{J-1}\frac{z^{2j+1}}{2j+1},\qquad
 U_J(z)=L_J(z)+\frac{2z^{2J+1}}{(2J+1)(1-z^2)}.
\]
With $J=24$, put
\[
 u_n=\left\lceil2^{40}\bigl(kU_J(1/3)+U_J(z)\bigr)\right\rceil.
\]
Formula \eqref{eq:log-series} proves the second inequality in
\eqref{eq:small-bound-certificate}.  The first is checked using only integers:
\begin{equation}\label{eq:certificate-integer-test}
 B_n(125u_n)^n<(99n\,2^{40})^n.
\end{equation}
The Bell recurrence generates $B_1,\ldots,B_{511}$.  The included verifier
performs precisely these 511 comparisons, without floating point arithmetic,
and all are true.  This is a finite computation proof: the arithmetic input,
all endpoints, and the complete algorithm are specified.  It can be expanded
into ground defining equations of \ipra, followed by their finite conjunction.
The Python run is independent reproducibility evidence, not a claim that a
Python interpreter checks the object logic.  Combining this finite certificate
with the already-proved infinite range proves the theorem.
\end{proof}

The near-maximal constant among these small indices occurs at $n=4$:
$15^{1/4}\log5/4$ is close to $0.792$.  This explains why ordinary
floating-point rounding is unsuitable as a proof of all the small cases.
The certificate uses an upper enclosure for each logarithm and an exact
strict integer comparison.

\begin{corollary}[An explicit version of the eventual page bound]
For every $n\ge2^{16}$,
\begin{equation}\label{eq:eventual-six-tenths}
 B_n<\left(\frac{e^{-3/5}n}{\log(n+1)}\right)^n.
\end{equation}
Consequently, for every positive rational $\eta$, the same threshold gives
$B_n<(e^{-3/5+\eta}n/\log(n+1))^n$.
\end{corollary}
\begin{proof}
Use $r=\tfrac34\log(n+1)$ in the coefficient bound.  The same cancellation
as before yields
\[
 \frac{B_n^{1/n}\log(n+1)}n
 \le\frac43\exp\left(-1+
             \frac{\log n+(n+1)^{3/4}}n\right).
\]
For $n\ge2^{16}$, the decreasing logarithm quotient is at most
$16/2^{16}$ because $\log2<1$.  Also
\[
 \frac{(n+1)^{3/4}}n
 =n^{-1/4}(1+1/n)^{3/4}
 \le\frac1{16}(1+2^{-16}).
\]
Their sum is less than $1/10$, by a rational comparison.
Since $e^{3/10}>1+3/10+(3/10)^2/2>4/3$, we have
$\log(4/3)<3/10$.  The displayed bound is therefore strictly less than
$\exp(-1+1/10+3/10)=e^{-3/5}$.
\end{proof}

The source page also supplies a particular threshold using an inverse
function $d^{-1}$.  The last corollary proves the same eventual inequality
with a different, explicit integer threshold and without computing that
inverse.  The sharper root-growth limit of Corollary~\ref{cor:root-growth}
further yields the optimal asymptotic constant $e^{-1}$: for every rational
$\eta>0$ there is a primitive recursively specified threshold beyond which
$B_n<(e^{-1+\eta}n/\log(n+1))^n$.
